% !TeX root = RJwrapper.tex
\title{Nowcast: Economic Nowcasting with Bridge Equations and Real-Time Evaluation in R}


\author{by Charles Coverdale}

\maketitle

\abstract{%
The nowcast package provides bridge-equation nowcasting in R. It estimates current-quarter values of a low-frequency macroeconomic variable (typically GDP) from higher-frequency indicators (retail sales, industrial production, payrolls, sentiment) that are released weeks or months ahead of the official target. Mixed-frequency alignment handles the ragged edge of real-time data. Pseudo-real-time backtesting over expanding or rolling windows produces an out-of-sample performance record, and the Diebold-Mariano test with the Harvey-Leybourne-Newbold finite-sample correction compares competing specifications on the same target. The package is pure R with cli as the only non-base import, is available on CRAN, and plugs into any source of quarterly or monthly data.
}

\section{Introduction}\label{introduction}

The \CRANpkg{nowcast} package implements the standard workflow for
bridge-equation nowcasting: mixed-frequency alignment, estimation
with optional autoregressive terms, and pseudo-real-time evaluation.
Ten exported functions share a uniform \code{nc\_} prefix and accept
plain data frames of dates and values, with no hard tie to
\CRANpkg{xts}, \CRANpkg{zoo}, or any data-vendor schema.

Nowcasting (the practice of estimating a macroeconomic variable
before it is officially released) has been standard in central banks
since the early 2000s. GDP is typically published six to eight weeks
after the quarter closes, but monthly indicators relevant to GDP
arrive much sooner: US retail sales with a two- to three-week lag,
industrial production with a two-week lag, nonfarm payrolls in the
first week of the following month. The appeal of a nowcast is that
it converts the flow of monthly releases during a quarter into a
continuously updated estimate of the target.

R has had partial support for this workflow. \CRANpkg{bridgr}
implements bridge equations within the tidyverse ecosystem.
\CRANpkg{midasr} covers the related MIDAS family of mixed-frequency
models. \CRANpkg{bigtime} fits large-scale time-series models
including sparse VARs at scale. None of these covers the full
evaluation pipeline that a nowcasting analyst needs:
pseudo-real-time backtesting with expanding or rolling windows, and
pairwise Diebold-Mariano testing with the Harvey-Leybourne-Newbold
finite-sample correction for small samples. \CRANpkg{nowcast} fills
that gap.

\section{Background}\label{background}

\textbf{The mixed-frequency problem.} A quarterly target (GDP) must be
aligned with monthly indicators. The standard choice, following
\citet{mariano2003new}, is to aggregate monthly indicators to a quarterly
frequency by averaging within the target quarter, which is the
bridge-equation formulation, or to treat the low-frequency variable
as a linear combination of latent high-frequency values, which is
the MIDAS and dynamic-factor approach.

\textbf{Bridge equations.} A bridge equation regresses the low-frequency
target on the quarterly averages of monthly indicators
\citep{baffigi2004bridge}. Optionally an autoregressive term in the
target's own lags accounts for persistence not captured by the
indicators. Bridge equations are simpler than MIDAS (which specifies
a high-frequency weighting function) and dynamic factor models
(which specify a latent state-space representation), but they remain
the workhorse at central banks because their inputs (the monthly
indicators) are the observable quantities that release-date
press-releases report. \citet{schumacher2016comparison} compares bridge to
MIDAS on a common dataset and finds comparable out-of-sample
accuracy with bridge's simpler specification; \citet{clements2008macroeconomic}
reach a similar conclusion for US GDP. \citet{foroni2013survey} provide a
survey of mixed-frequency econometrics covering both approaches.

\textbf{The ragged edge.} Different indicators have different publication
lags. On any given calendar date, retail sales may be published
through month \(t-1\), industrial production through \(t-2\), employment
through \(t-1\), and consumer sentiment through \(t\) flash estimate.
The resulting data matrix has a jagged right edge; a nowcasting
model must handle this. \citet{wallis1986forecasting} formalised the
ragged-edge problem in the context of mixed-frequency econometrics,
and all modern nowcasting implementations address it explicitly.

\textbf{Real-time evaluation.} In-sample fit is a weak guide to nowcast
quality because it ignores both the ragged-edge publication pattern
and out-of-sample instability. The canonical alternative is
pseudo-real-time evaluation \citep{stark2002database}: for each target
date, estimate the model on data available at that date only, make
a nowcast, compare to the eventual realisation. An expanding window
re-estimates at each step on all prior data; a rolling window fixes
the estimation sample length. \citet{giannone2008nowcasting} standardised
this approach. \citet{aruoba2009real} construct a related continuous-time
index of US business conditions that has become the canonical
dynamic-factor comparator for any single-equation nowcast.
\citet{bok2018macroeconomic} survey the modern nowcasting literature,
including the New York Fed Nowcast (the institutional benchmark for
US GDP nowcasting).

\textbf{Comparing nowcasts.} The Diebold-Mariano test
\citep{diebold1995comparing} compares loss differentials between two
forecast series against the null of equal predictive accuracy.
\citet{harvey1997testing} showed that the original test is over-sized in
small samples and proposed a finite-sample correction; the
HLN-corrected statistic is the version implemented by
\code{nc\_dm\_test()}.

\section{Package design}\label{package-design}

\subsection{Architecture}\label{architecture}

\CRANpkg{nowcast} is pure R with no compiled code. Runtime imports
are \CRANpkg{cli}, \code{grDevices}, \code{graphics}, and
\code{stats}. R 4.1.0 or later is required. The test suite
contains over one hundred and fifty tests covering every exported
function.

\subsection{Uniform function interface}\label{uniform-function-interface}

Every exported function is prefixed \code{nc\_}. Data functions
(\code{nc\_align}, \code{nc\_ragged\_edge}, \code{nc\_aggregate},
\code{nc\_transform}) accept or return data frames with a
\code{date} column and one or more value columns. Modelling
functions (\code{nc\_bridge}, \code{nc\_backtest},
\code{nc\_evaluate}, \code{nc\_dm\_test}) accept formula syntax
(\code{target \~{} ind1 + ind2}) and return S3 objects.

\subsection{S3 classes and methods}\label{s3-classes-and-methods}

Alignment produces an \code{nc\_dataset} object carrying the
quarterly-aggregated data plus ragged-edge metadata. Modelling
returns \code{nowcast\_result}, \code{nowcast\_backtest}, or
\code{nowcast\_dm} objects, each with \code{print()},
\code{summary()}, and where appropriate \code{plot()} methods.
\code{predict()} on a fitted \code{nowcast\_result} scores new data
without re-estimating.

\subsection{Reproducibility}\label{reproducibility}

Every numerical result is deterministic given the input. No random
initialisation, no stochastic optimiser, no cached vintage state.
The package ships no bundled data; all examples in this paper use US
series obtained from the Federal Reserve Bank of St.~Louis Economic
Data (FRED) service.

\section{Alignment and the ragged edge}\label{alignment-and-the-ragged-edge}

\code{nc\_align()} takes a quarterly target data frame and one or
more monthly indicator data frames, aggregates each monthly series
to quarterly frequency by averaging within the target quarter, and
returns an \code{nc\_dataset} with the joined data plus diagnostic
columns recording how many monthly observations contributed to each
quarter. Partial quarters (where the monthly indicator has fewer
than three months within the target) are flagged rather than
silently dropped.

\code{nc\_ragged\_edge()} summarises the right edge of the
alignment: for each indicator, how far past the latest target
observation is the indicator's own latest observation? Figure
\ref{fig:ragged} shows the pattern for four US monthly indicators
over the most recent twenty-four months.

\begin{figure}

{\centering \includegraphics[width=1\linewidth,alt={A heatmap with four rows (retail sales, industrial production, nonfarm payrolls, consumer sentiment) and 24 columns (months). Blue tiles indicate months with published observations; grey tiles indicate months not yet released. The right edge is jagged: sentiment extends furthest, payrolls and retail sales end one month earlier, industrial production two months earlier. This illustrates the ragged-edge publication pattern.}]{figures/fig1_ragged_edge} 

}

\caption{Availability heatmap of four US monthly indicators over the last 24 months. Blue tiles mark months with published observations, grey tiles mark months where the series has not yet been released. The jagged right edge at the most recent dates is the problem \texttt{nc\_align()} solves: retail sales, industrial production, payrolls, and sentiment all have different publication lags.}\label{fig:ragged}
\end{figure}

\code{nc\_aggregate()} and \code{nc\_transform()} expose the
aggregation and stationarity-transformation steps separately for
users who need manual control: log-difference, first-difference,
percent change, and year-over-year transforms are all supported.

\section{Bridge equations}\label{bridge-equations}

\code{nc\_bridge()} fits the bridge equation by ordinary least
squares, with an optional autoregressive term in the target
(\code{ar\_order = 1} by default). The return object carries the
point nowcast, a standard error, a confidence interval at the
user-specified level, and the underlying fitted \code{lm} model for
users who want coefficients, fitted values, and diagnostics. The
\code{newdata} argument allows nowcasting out of sample without
re-fitting.

Figure \ref{fig:backtest} shows a pseudo-real-time backtest of a
bridge equation with four US monthly indicators (retail sales,
industrial production, nonfarm payrolls, and consumer sentiment) as
regressors, targeting US real GDP growth over 2012 to present. The
nowcast tracks the target through two episodes of pronounced
volatility: the 2020 pandemic contraction and the post-2022 recovery.

\begin{figure}

{\centering \includegraphics[width=1\linewidth,alt={A time-series line chart from 2012 to the mid-2020s. A grey solid line shows realised quarterly US real GDP log-differences; a blue dashed line shows the one-step-ahead bridge-equation nowcast. Both series move together through the period, with the steepest downward spike in 2020 followed by a symmetric recovery. The nowcast closely tracks the realised series outside of the pandemic quarters.}]{figures/fig2_backtest} 

}

\caption{Pseudo-real-time bridge-equation nowcast of quarterly US real GDP growth, 2012 onwards. Grey solid line is the realised log-difference of real GDP from FRED series \texttt{GDPC1}. Blue dashed line is the one-step-ahead nowcast from \texttt{nc\_backtest()} with four monthly indicators aggregated to quarterly frequency. Nowcast RMSE over the evaluation window is 0.95 percentage points, MAE is 0.51, bias is 0.03.}\label{fig:backtest}
\end{figure}

Figure \ref{fig:coeffs} shows the estimated coefficients and
confidence intervals for the four indicators in the full-sample
bridge equation, interpreted as the response of quarterly GDP growth
to a one-percentage-point increase in the quarterly average of each
monthly indicator. Retail sales and industrial production have the
expected positive sign and the tightest intervals; payrolls and
sentiment contribute more modestly.

\begin{figure}

{\centering \includegraphics[width=1\linewidth,alt={A horizontal dot-whisker plot showing four coefficient estimates for retail sales, industrial production, nonfarm payrolls, and consumer sentiment. Retail sales and industrial production have the largest positive point estimates and tightest confidence intervals. Payrolls and sentiment have smaller estimates with wider intervals that straddle zero. A vertical reference line at zero is shown.}]{figures/fig4_coefficients} 

}

\caption{Bridge-equation coefficients for four monthly indicators, full-sample fit. Points are ordinary-least-squares estimates; horizontal bars are 95 per cent confidence intervals. Vertical line at zero marks no effect. Retail sales and industrial production dominate; payrolls and sentiment contribute marginally in this specification.}\label{fig:coeffs}
\end{figure}

\section{Pseudo-real-time evaluation}\label{pseudo-real-time-evaluation}

\code{nc\_backtest()} walks through the sample one target at a time.
At each step it fits the bridge equation on data available through
the previous quarter and makes a nowcast for the next. The
\code{window} argument selects an expanding window (all prior data)
or a rolling window of fixed length, following
\citet{giannone2008nowcasting}. The \code{start} argument sets the first
evaluation point; everything prior is used for initial estimation.

Figure \ref{fig:window} compares expanding and rolling-window
errors on the same US GDP target. The expanding window benefits
from all available data and produces a marginally smaller RMSE; the
rolling window (forty-quarter fixed length) is more responsive to
structural change, at the cost of ignoring older observations.

\begin{figure}

{\centering \includegraphics[width=1\linewidth,alt={A time-series line chart of nowcast errors over the evaluation window. A blue solid line is the expanding-window error; a red dashed line is the rolling-window error. Both trace near zero through most of the sample with a large positive spike around 2020. The rolling window returns toward zero more quickly after the spike than the expanding window.}]{figures/fig3_window} 

}

\caption{Nowcast error (nowcast minus realised) from expanding-window and rolling-window backtests of the same bridge-equation specification. Expanding window (blue, solid) uses all prior data at each step; rolling window (red, dashed) uses the preceding 40 quarters only. Both peak around the 2020 pandemic contraction; the rolling specification is more responsive afterwards.}\label{fig:window}
\end{figure}

\code{nc\_evaluate()} takes a realised series and a nowcast series
and returns root-mean-squared error, mean absolute error, and bias.
These are the standard metrics reported alongside any published
nowcasting exercise.

\section{Testing equal predictive accuracy}\label{testing-equal-predictive-accuracy}

\code{nc\_dm\_test()} implements the Diebold-Mariano test for equal
predictive accuracy between two forecast series on the same target.
It accepts two loss vectors (\code{e1}, \code{e2}) and returns the
test statistic, \(p\)-value, and interpretation. The \code{h} argument
sets the forecast horizon (one for a same-quarter nowcast). The
\code{loss} argument selects squared or absolute loss.
\citet{harvey1997testing} showed that the asymptotic test is over-sized in
samples smaller than about fifty; the HLN correction applied here
scales the statistic by a factor depending on sample size and
horizon to restore approximate nominal coverage.

Figure \ref{fig:dm} compares two specifications of the bridge
equation on the same US GDP target: with and without the AR(1) term.
The RMSEs are nearly identical and the DM test fails to reject the
null of equal predictive accuracy. The AR term is neither helpful
nor harmful in this example, consistent with the modest persistence
of quarterly log-differenced GDP growth.

\begin{figure}

{\centering \includegraphics[width=1\linewidth,alt={A two-bar chart comparing the RMSE of two bridge-equation specifications on US GDP. The left bar is the specification with an AR(1) term; the right bar is the specification without. Both bars are nearly identical in height, near 0.95 percentage points. An annotation box reports the HLN-corrected DM statistic and a p-value above conventional significance thresholds.}]{figures/fig5_dm} 

}

\caption{Diebold-Mariano comparison of two bridge-equation specifications on US GDP. Bars show backtest RMSE. Left: bridge with an AR(1) term. Right: bridge without. The HLN-corrected Diebold-Mariano test statistic and p-value are reported in the annotation box. The test fails to reject the null of equal predictive accuracy.}\label{fig:dm}
\end{figure}

\section{Replication}\label{replication}

The canonical workflow is four lines.

\begin{verbatim}
aligned  <- nc_align(gdp, retail = retail_df, indpro = indpro_df)
result   <- nc_bridge(target ~ retail + indpro, data = aligned)
backtest <- nc_backtest(target ~ retail + indpro, data = aligned, start = 40)
backtest$metrics
\end{verbatim}

Line one aligns the quarterly target with the monthly indicators,
averaging each indicator within its target quarter. Line two fits
the bridge equation. Line three walks a pseudo-real-time expanding
window through the sample. Line four pulls out the RMSE, MAE, and
bias.

Table \ref{tab:metrics} compares three bridge-equation
specifications against a naive AR(1) benchmark that uses only the
target's own lag. The full four-indicator bridge delivers an RMSE
of 0.95 percentage points, roughly 46 per cent below the naive
benchmark's 1.75. Dropping the AR term adds about three per cent to
RMSE. Dropping all indicators except retail sales and industrial
production nearly doubles the RMSE back to the naive-benchmark
level, confirming that the payrolls and sentiment series carry
meaningful marginal signal. The pairwise Diebold-Mariano test
between the full bridge and the naive benchmark returns \(p = 0.17\),
indicating that the sample is not quite large enough to reject
equal accuracy at conventional levels, even though the point
estimate strongly favours the bridge.

\begin{table}

\caption{\label{tab:metrics}Backtest metrics for three bridge-equation specifications on US real GDP growth. Expanding window, first evaluation at quarter forty, 2012-present evaluation sample. Adding the AR(1) term marginally reduces RMSE; dropping two of the four indicators modestly widens it.}
\centering
\begin{tabular}[t]{lrrr}
\toprule
Specification & RMSE & MAE & Bias\\
\midrule
Naive AR(1) on target only (benchmark) & 1.751 & 0.734 & -0.148\\
Bridge + AR(1), all four indicators & 0.950 & 0.509 & 0.025\\
Bridge no AR, all four indicators & 0.974 & 0.501 & 0.081\\
Bridge + AR(1), retail + INDPRO only & 1.827 & 0.743 & -0.201\\
\bottomrule
\end{tabular}
\end{table}

\section{A case study of the 2020 pandemic quarter by quarter}\label{a-case-study-of-the-2020-pandemic-quarter-by-quarter}

The expanding-window backtest provides a natural narrative window
for the 2020 pandemic shock, the largest peacetime real-GDP
contraction in US history. Evaluated quarter by quarter, the bridge
nowcast returned:

\begin{itemize}
\tightlist
\item
  2020 Q1: nowcast \(-1.4\) per cent (log-difference), realised
  \(-1.3\) per cent. The four monthly indicators captured the
  front-end hit nearly exactly.
\item
  2020 Q2: nowcast \(-3.2\) per cent, realised \(-8.2\) per cent (a
  \(-31\) per cent annualised print). The nowcast caught the
  direction but missed the magnitude by five percentage points.
  Retail sales, industrial production, and payrolls all fell
  sharply, but the simultaneous collapse in services consumption,
  which is not directly captured by any of the four indicators in
  this specification, was unprecedented.
\item
  2020 Q3: nowcast \(+4.7\) per cent, realised \(+7.5\) per cent. Again
  the direction was captured but the magnitude understated, for the
  same reason in reverse: the services rebound exceeded what any
  linear combination of the four monthly series would have
  predicted.
\item
  2020 Q4: nowcast \(-1.6\) per cent, realised \(+1.1\) per cent. By
  the end of the year the model had absorbed the new-regime signal
  from the 2020 Q3 rebound and began over-correcting.
\end{itemize}

The episode illustrates both the appeal and the limit of
bridge-equation nowcasts. The model is easy to estimate, reports
sensible uncertainty, and captures direction reliably; it cannot
capture shocks that move the high-frequency target through a
channel not in the indicator set. This is the common pattern: add
one or two services or spending indicators (personal consumption
expenditure on services, for example) and the magnitude misses
narrow substantially.

\section{Limitations}\label{limitations}

Five limitations apply.

\begin{enumerate}
\def\labelenumi{\arabic{enumi}.}
\tightlist
\item
  \CRANpkg{nowcast} implements bridge equations only. Mixed-data
  sampling (MIDAS) regressions of the \citet{ghysels2004midas} family, and
  dynamic factor models in the \citet{giannone2008nowcasting} tradition,
  are out of scope; users needing those should look to
  \CRANpkg{midasr} and \CRANpkg{bigtime} respectively.
\item
  The package does not fetch data. Users supply a data frame of
  dates and values from FRED, Eurostat, ONS, or any other source.
  \CRANpkg{fredr} is a natural partner for US data;
  \CRANpkg{readecb} for Euro-area data; \CRANpkg{ons} for UK.
\item
  Real-time data revisions are not tracked. The package operates
  on final-vintage data by default; true real-time nowcast
  evaluation requires vintage-aware data frames which the user must
  construct (for US data, ALFRED at FRED supplies them).
\item
  Flash-estimate indicators with intra-month release dates require
  manual alignment. The monthly aggregation in \code{nc\_align()}
  assumes one observation per month.
\item
  The AR specification in \code{nc\_bridge()} is restricted to
  non-negative integer orders of the target's own lags.
  Autoregressive-moving-average or fractional specifications are
  out of scope.
\end{enumerate}

\section{Appendix of formula definitions}\label{appendix-of-formula-definitions}

\textbf{Bridge equation with AR(\(p\)) term.} Let \(y_t\) denote the
quarterly target and \(x_{j,t}^{\mathrm{agg}}\) the quarterly
aggregation of monthly indicator \(j\). The bridge specification is
\[y_t = \alpha + \sum_{j=1}^{J} \beta_j x_{j,t}^{\mathrm{agg}}
  + \sum_{k=1}^{p} \rho_k y_{t-k} + \varepsilon_t,\]
estimated by ordinary least squares. Nowcasts substitute the
contemporaneous quarterly aggregate of each indicator; if the
indicator's final month is missing, the aggregate uses the
partial-quarter average, flagged in the package's ragged-edge
metadata.

\textbf{Mincer-Zarnowitz regression.} For realised \(y_t\) and forecast
\(\hat{y}_t\), estimate \(y_t = \alpha + \beta \hat{y}_t + \varepsilon_t\)
and test \(H_0: (\alpha, \beta) = (0, 1)\). Rejection indicates bias
(\(\alpha \neq 0\)) or systematic under- or over-reaction
(\(\beta \neq 1\)).

\textbf{Diebold-Mariano test with HLN correction.} Let
\(d_t = L(e_{1,t}) - L(e_{2,t})\) be the loss differential between
two forecast errors under a chosen loss function \(L\). The
Diebold-Mariano statistic is \(\mathrm{DM} = \bar{d} / \sqrt{V_n}\),
where \(V_n\) is a heteroscedasticity-and-autocorrelation-consistent
variance estimate. \citet{harvey1997testing} show this is over-sized in
small samples and propose the correction
\[\mathrm{DM}^* = \mathrm{DM} \cdot
  \sqrt{\frac{n + 1 - 2h + h(h-1)/n}{n}},\]
with \(p\)-values evaluated against a \(t_{n-1}\) distribution. This is
the statistic \code{nc\_dm\_test()} returns.

\section{Conclusion}\label{conclusion}

Nowcasting is a core activity at every central bank's research
department and a frequent task in sovereign fiscal offices and
macro-hedge funds, and it had been only partially served on CRAN.
\CRANpkg{nowcast} closes the gap with a complete bridge-equation
workflow: mixed-frequency alignment, estimation with optional
autoregressive terms, pseudo-real-time backtesting over expanding or
rolling windows, and HLN-corrected Diebold-Mariano testing.
Every function operates on plain data frames, has no hard
dependency on data-vendor packages, and returns S3 objects with
print, summary, and plot methods. Planned additions include a
mixed-data-sampling interface, dynamic-factor estimation for large
indicator panels, real-time vintage support, and flash-estimate
handling.

\section*{Acknowledgements}\label{acknowledgements}
\addcontentsline{toc}{section}{Acknowledgements}

The Federal Reserve Bank of St.~Louis publishes the FRED service
used for every example in this paper. The CRAN team maintain the
infrastructure that makes R package distribution possible.

\section*{References}
\bibliography{RJreferences.bib}

\address{%
Charles Coverdale\\
Independent\\%
London, United Kingdom\\
%
\url{https://github.com/charlescoverdale/nowcast}\\%
%
\href{mailto:charles.f.coverdale@gmail.com}{\nolinkurl{charles.f.coverdale@gmail.com}}%
}
